
        \documentclass[spanish, 11pt]{exam}

        %These tell TeX which packages to use.
        \usepackage{array,epsfig}
        \usepackage{amsmath, textcomp}
        \usepackage{amsfonts}
        \usepackage{amssymb}
        \usepackage{amsxtra}
        \usepackage{amsthm}
        \usepackage{mathrsfs}
        \usepackage{color}
        \usepackage{multicol, xparse}
        \usepackage{verbatim}


        \usepackage[utf8]{inputenc}
        \usepackage[spanish]{babel}
        \usepackage{eurosym}

        \usepackage{graphicx}
        \graphicspath{{../img/}}



        \printanswers
        \nopointsinmargin
        \pointformat{}

        %Pagination stuff.
        %\setlength{\topmargin}{-.3 in}
        %\setlength{\oddsidemargin}{0in}
        %\setlength{\evensidemargin}{0in}
        %\setlength{\textheight}{9.in}
        %\setlength{\textwidth}{6.5in}
        %\pagestyle{empty}

        \let\multicolmulticols\multicols
        \let\endmulticolmulticols\endmulticols
        \RenewDocumentEnvironment{multicols}{mO{}}
         {%
          \ifnum#1=1
            #2%
          \else % More than 1 column
            \multicolmulticols{#1}[#2]
          \fi
         }
         {%
          \ifnum#1=1
          \else % More than 1 column
            \endmulticolmulticols
          \fi
         }
        \renewcommand{\solutiontitle}{\noindent\textbf{Sol:}\enspace}

        \newcommand{\samedir}{\mathbin{\!/\mkern-5mu/\!}}

        \newcommand{\class}{1º Bachillerato}
        \newcommand{\examdate}{\today}

        \newcommand{\tipo}{A}


        \newcommand{\timelimit}{50 minutos}



        \pagestyle{head}
        \firstpageheader{\includegraphics[width=0.2\columnwidth]{header_left}}{\textbf{Departamento de Matemáticas\linebreak \class}\linebreak \examnum}{\includegraphics[width=0.1\columnwidth]{header_right}}
        \runningheader{\class}{\examnum}{Página \thepage\ of \numpages}
        \runningheadrule

        \newcommand{\examnum}{6 - Ecuaciones de segundo grado}
        \begin{document}
        \begin{questions}
        \question p016e01 - Resuelve las ecuaciones:
        \begin{multicols}{2} 
        \begin{parts} \part[1]  $ x^{2} + 6 = 0 $  \begin{solution}  $ \emptyset $  \end{solution} \part[1]  $ x^{2} - 9 = 0 $  \begin{solution}  $ \left\{-3, 3\right\} $  \end{solution} \part[1]  $ x^{2} + 3 x = 0 $  \begin{solution}  $ \left\{-3, 0\right\} $  \end{solution} \part[1]  $ 3 x^{2} - 11 x = 0 $  \begin{solution}  $ \left\{0, \frac{11}{3}\right\} $  \end{solution} \part[1]  $ 4 x^{2} - 32 x = 0 $  \begin{solution}  $ \left\{0, 8\right\} $  \end{solution} \part[1]  $ 5 x^{2} = 0 $  \begin{solution}  $ \left\{0\right\} $  \end{solution} \part[1]  $ 12 x^{2} - 18 = 0 $  \begin{solution}  $ \left\{- \frac{\sqrt{6}}{2}, \frac{\sqrt{6}}{2}\right\} $  \end{solution} \part[1]  $ 3 \left(- x + 1\right) \left(x + 1\right) = 3 $  \begin{solution}  $ \left\{0\right\} $  \end{solution} \part[1]  $ 3 \left(x^{2} - 2\right) = 21 $  \begin{solution}  $ \left\{-3, 3\right\} $  \end{solution}
        \end{parts}
        \end{multicols}
        \question p016e02 - Resuelve las ecuaciones:
        \begin{multicols}{2} 
        \begin{parts} \part[1]  $ \left(2 x^{2} + 11 x\right) - 6 = 0 $  \begin{solution}  $ \left\{-6, \frac{1}{2}\right\} $  \end{solution} \part[1]  $ \left(x^{2} - 10 x\right) + 25 = 0 $  \begin{solution}  $ \left\{5\right\} $  \end{solution} \part[1]  $ \left(x^{2} + x\right) + 1 = 0 $  \begin{solution}  $ \emptyset $  \end{solution} \part[1]  $ \left(x^{2} - 2 x\right) - 1 = 0 $  \begin{solution}  $ \left\{1 + \sqrt{2}, - \sqrt{2} + 1\right\} $  \end{solution} \part[1]  $ \left(3 x^{2} + 5 x\right) - 2 = 0 $  \begin{solution}  $ \left\{-2, \frac{1}{3}\right\} $  \end{solution} \part[1]  $ \left(4 x^{2} - 4 x\right) + 1 = 0 $  \begin{solution}  $ \left\{\frac{1}{2}\right\} $  \end{solution} \part[1]  $ \left(2 x^{2} - 9 x\right) + 11 = 0 $  \begin{solution}  $ \emptyset $  \end{solution}
        \end{parts}
        \end{multicols}
        \question p016e03 - Resuelve las ecuaciones:
        \begin{multicols}{2} 
        \begin{parts} \part[1]  $ - x \left(x - 2\right) + 9 = 4 x + 6 $  \begin{solution}  $ \left\{-3, 1\right\} $  \end{solution} \part[1]  $ - \left(x - 1\right) \left(x + 4\right) + 2 \left(x^{2} - 3\right) = x - 2 $  \begin{solution}  $ \left\{0, 4\right\} $  \end{solution} \part[1]  $ x \left(x - 1\right) - 2 \left(x - 3\right) \left(x - 2\right) = 2 $  \begin{solution}  $ \left\{2, 7\right\} $  \end{solution} \part[1]  $ \left(2 x^{2} - 11 x\right) + 12 = 0 $  \begin{solution}  $ \left\{\frac{3}{2}, 4\right\} $  \end{solution} \part[1]  $ 3 \left(x - 1\right) \left(x + 2\right) = 0 $  \begin{solution}  $ \left\{-2, 1\right\} $  \end{solution} \part[1]  $ \left(x - 2\right)^{2} = 3 $  \begin{solution}  $ \left\{- \sqrt{3} + 2, \sqrt{3} + 2\right\} $  \end{solution} \part[1]  $ 21 x - 100 = - x + \left(x^{2} + 21\right) $  \begin{solution}  $ \left\{11\right\} $  \end{solution} \part[1]  $ \frac{x}{3} \left(x - \frac{1}{6}\right) = x - 1 $  \begin{solution}  $ \emptyset $  \end{solution} \part[1]  $ \left(- \frac{x}{3} + \frac{1}{3}\right) + \frac{2 x^{2} + 1}{2} = \frac{- x + 1}{6} $  \begin{solution}  $ \emptyset $  \end{solution}
        \end{parts}
        \end{multicols}
        
    \end{questions}
    \end{document}
    